\documentclass[a4paper]{article}

\usepackage{graphicx}
\usepackage{amsfonts}
\usepackage{geometry}
\usepackage[utf8]{inputenc}
\usepackage{array}
\usepackage{calc}
\usepackage{layout}
\geometry{
	a4paper,
	total={170mm,257mm},
	left=15mm,
	top=15mm,
}


\title{\begin{center}
		\setlength{\tabcolsep}{0pt}
		\begin{tabular}{>{\raggedleft}m{3cm}>{\centering}m{\textwidth - 5cm\relax}>{\raggedright}m{3cm}}
			\includegraphics[width=\linewidth]{poliba.png}%
			&%
			\textbf{ } \\[5pt]%
			\textbf{\ }%
			&%
			\includegraphics[width=\linewidth]{dei.png} %
		\end{tabular}
	\end{center}
	\textbf{Relazione attività di laboratorio\\}
	\textbf{\large Esercitazione 1\\- Corso di Fondamenti delle Misurazione -\\ }
	{\normalsize 
		- Ingegneria Informatica e dell'Automazione -
}}
\author{A cura di:\\
	Laura Loperfino\\
	Andrea Lops \\
	Teresa Pantone\\
	Paolo Rotolo\\
	\\
	Prof. Ing. Filippo Attivissimo
}

\date{29/11/2018}


\begin{document}

\maketitle

\section{Introduzione}

Lo scopo di questa esercitazione consiste nell'effettuare differenti misure, dirette e non, con gli strumenti tipici di un laboratorio quale multimetro e oscilloscopio,  interpretando in modo corretto i valori dati dal datasheet di ogni strumento per estrapolare le informazioni necessarie ad esprimere queste in modo coerente e corretto.

\section{Calcoli}

Per poter calcolare le incertezze richieste abbiamo usato le seguenti \emph{formule}:\\
\subsection {Incertezza sulla resistenza \emph{R}}

In primo luogo, è stata effettuata la misura della resistenza tramite il multimetro \emph{Agilent 34401}:

\begin{figure}[htp]
	\centering
	\includegraphics[scale=0.20]{"figure1".jpg}
	\caption{Multimetro in misura}
	\label{}
\end{figure}
\noindent 
Per il calcolo dell'incertezza relativa alla resistenza si consulta la seguente tabella (\emph{Figure 2}) delle specifiche relative al multimetro \emph{Agilent 34401}:
\begin{center}
	\begin{figure}[htp]
		\centering
		\includegraphics[scale=0.50]{"figure2".png}
		\caption{Tabella di accuratezza del multimetro}
		\label{a}
	\end{figure}
\end{center}

Il valore della resistenza letto sul multimetro è  9.8887 k$\Omega$. Si sceglie il range appropriato per la misura (10 kOhm) e si leggono i valori di incertezza di lettura e di fondo scala.

\begin{Large} 
	\begin{equation} 
		U_{R} = [\pm 0.010\%rdg \pm 0.001\% FSO]
	 \end{equation}
	 
	 \begin{eqnarray*} 
	 	R = 9.8887 \pm 0.0010 k\Omega 
	 \end{eqnarray*}
\end{Large}

\subsection {Incertezza sui tempi}

Per calcolare l'incertezza sulle misure è necessario conoscere le specifiche di performance dell'\emph{Oscilloscopio HP54603B}. Si prende come riferimento il sistema orizzontale, ovvero quello riferito al tempo, avendo appurato che tali misure si riferiscano ad esso.

\begin{figure}[htp]
	\centering
	\includegraphics[scale=0.20]{"figure3".jpg}
	\caption{Oscilloscopio}
	\label{}
\end{figure}

Di seguito sono riportati i calcoli relativi all'incertezza che verranno illustrati solo per la prima misura in quanto sono gli stessi che verranno utilizzati per le misure successive.

\subsubsection{Tempo di salita con sonda non compensata}

Per calcolare l'incertezza in questo caso si fa riferimento al valore \emph{cursor accuracy} presente nel datasheet dell'oscilloscopio (\emph{Figure 4}), riferito all'asse dei tempi:

\begin{figure}[htp]
	\centering
	\includegraphics[scale=0.50]{"figplus".jpg}
	\caption{Datasheet oscilloscopio HP54603B}
	\label{}
\end{figure}


\begin{Large}
	\begin{equation}
  		U_{t}=[\pm 0.01\%rdg\pm 0.2\%FSO \pm 200ps]
	\end{equation}
\end{Large}\\

Si noti che per poter calcolare l'incertezza si necessita del valore di fondo scala ottenuto moltiplicando l'impostazione di guadagno per il numero di divisioni.

Nel caso del tempo di salita con la sonda non compensata, l'impostazione di guadagno verticale/base dei tempi è 10.0V/2.0$\mu$s.
Il numero di divisioni è 10, quindi:
\begin{Large}
	\begin{equation}
  		FS=2.0\mu s * 10= 20\mu s
	\end{equation}
\end{Large}\\

Quindi si ricava l'incertezza:
\begin{Large}
	\begin{equation}
  		{t_r}= 5.50\mu s\pm (+ 0,00055\mu s + 0,04\mu s + 200ps)
	\end{equation}
\end{Large}\\
Per rendere più elegante la scrittura del risultato senza intaccarne la validità, si decide di approssimarlo alla seconda cifra decimale.
\begin{Large}
	\begin{equation}
  		{t_r}= 5.50\pm 0,04 \mu s
	\end{equation}
\end{Large}\\
\subsubsection{Tempo di discesa con sonda non compensata}
\begin{Large}
	\begin{equation}
  		{t_f}= 5.70\mu s\pm (+ 0,00057\mu s + 0,04\mu s + 200ps)
	\end{equation}
\end{Large}\\

Arrotondando: 

\begin{Large}
	\begin{equation}
  		{t_f}= 5.70\pm 0,04 \mu s
	\end{equation}
\end{Large}\\

Una sonda compensata è composta da una resistenza elevata in parallelo ad un condensatore di capacità variabile (compensatore).

Grazie a quest'ultima, è possibile annullare l'effetto di carico per segnali ad alte frequenze (più di 5 kHz) e di ampiezza sufficientemente alta rispetto al rumore.

\subsubsection{Tempo di salita con sonda compensata 10x}

\begin{Large}
	\begin{equation}
		{t_{r_{10x}}}= 975ns \pm( + 0,0975ns + 10ns + 200ps)
	\end{equation}
\end{Large}\\

Arrotondando: 

\begin{Large}
	\begin{equation}
		{t_{r_{10x}}}= 975\pm 10 ns
	\end{equation}
\end{Large}\\

\subsubsection{Tempo di salita con sonda non compensata e inserzione di una resistenza}

\begin{Large}
	\begin{equation}
		t_{r_R}= 7.11\mu s \pm( + 0,00071\mu s + 0,04\mu s + 200ps)
	\end{equation}
\end{Large}\\

Arrotondando: 

\begin{Large}
	\begin{equation}
		t_{r_R}= 7.11\pm 0,04 \mu s
	\end{equation}
\end{Large}\\

\subsubsection{Tempo di salita con sonda compensata 10x e inserzione di una resistenza}

\begin{Large}
	\begin{equation}
		{t_{r_{R_{10x}}}}= 1300ns \pm ( + 0,13ns + 10ns + 200ps)
	\end{equation}
\end{Large}\\

Arrotondando: 

\begin{Large}
	\begin{equation}
		{t_{r_{R_{10x}}}}= 1300\pm 10 ns
	\end{equation}
\end{Large}\\

\subsubsection{Periodo segnale sinusoidale}
\begin{Large}
	\begin{equation}
  		{T_s}= 1000\mu s\pm (+ 0.1 \mu s + 4 \mu s + 200ps)
	\end{equation}
\end{Large}\\
\\
\\
\\

Arrotondando: 

\begin{Large}
	\begin{equation}
		{T_s}= 1000 \pm 4\mu s 
	\end{equation}
\end{Large}\\


\subsubsection{Frequenza segnale sinusoidale}
%Essendo una misura indiretta l'incertezza totale è data dalla somma delle derivate parziali delle %componenti:
%
%\begin{Large}
%	\begin{equation}
%	{U_{f_s}}= |\frac{d\frac{1}{T_s}}{d{T_s}}|U_{{T_s}}
%	\end{equation}
%\end{Large}\\
%Quindi l'incertezza:
%\begin{Large}
%	\begin{equation}
%	{U_{f_s}}= \frac{1}{1000000}*4.1=4.1 Hz
%	\end{equation}
%\end{Large}\\
Per quanto riguarda l'incertezza della frequenza, il calcolo è banale in quanto si ricorre alla stessa incertezza relativa al periodo del segnale sinusoidale con diversa unità di misura(\emph{Hz})
\begin{Large}
	\begin{equation}
  		{f_s}= 1000\pm 4Hz
	\end{equation}
\end{Large}\\

\subsection{Capacità sistema sonda non compensata, resistenza e oscilloscopio}
Diversamente dalle altre misure, il calcolo della capacità è una misura indiretta, essendo: 
\begin{Large}
	\begin{equation}
  		C= \frac{t_{r_R}}{R*ln(9)}
	\end{equation}
\end{Large}\\

Quindi:
\begin{Large}
	\begin{equation}
  		C= \frac{7.1}{9.9*ln(9)} [\frac{\mu s}{k\Omega}]= 0.325nF
	\end{equation}
\end{Large}\\

Essendo una misura indiretta l'incertezza totale è data dalla somma delle derivate parziali delle componenti:
\begin{Large}
	\begin{equation}
		U_{tot}=\sum_{i=1}^{N}|{\frac{df}{dy_i}}| U_i
	\end{equation}
\end{Large}\\

Quindi l'incertezza:
\begin{Large}
	\begin{equation}
		{U_C}= \frac{d\frac{t_{r_R}}{R*ln(9)}}{dt_{r_R}}U_{t_{r_R}}+\frac{d\frac{t_{r_R}}{R*ln(9)}}{dR}U_{R}
	\end{equation}
\end{Large}\\

Quindi:
\begin{Large}
	\begin{equation}
		{U_C}= 0.046*0.041+\frac{3.227}{9.9^2}*0.001=0.002 nF
	\end{equation}
\end{Large}\\

In conclusione: 
\begin{Large}
	\begin{equation}
		C=0.325\pm 0.002nF
	\end{equation}
\end{Large}\\

\subsubsection{Capacità sistema sonda compensata 10x, resistenza e oscilloscopio }
Si evita la rindondanza dei calcoli presenti precedentemente e si provvede ad eseguire il calcolo
\begin{Large}
	\begin{equation}
  		C_{10x}= \frac{1300}{9.9*ln(9)} [\frac{ns}{k\Omega}]= 59.688pF
	\end{equation}
\end{Large}\\

Quindi l'incertezza:
\begin{Large}
	\begin{equation}
		{U_{C_{10x}}}= 0.046*10.13+\frac{590.909}{9.9^2}*0.001=0.472 pF
	\end{equation}
\end{Large}\\

Quindi in conclusione:
\begin{Large}
	\begin{equation}
		{C_{10x}}= 59.7 \pm0.5 pF
	\end{equation}
\end{Large}\\

\subsection{Tensione picco-picco segnale sinusoidale} 
Per quanto riguarda la misura della tensione picco-picco del segnale sinusoidale, si si fa riferimento la \emph{dual cursor accuracy} (\emph{Figure 4}) :

\begin{Large}
	\begin{equation}
		U_{V_{p-p_{s}}}= [\pm 2.4\% rdg\pm 0.4\%FSO]
	\end{equation}
\end{Large}\\

Quindi: 
\begin{Large}
	\begin{equation}
		V_{p-p_{s}}= 3.938V\pm 0.0945V\pm 0.032V
	\end{equation}
\end{Large}\\

Arrotondando:
\begin{Large}
	\begin{equation}
		V_{p-p_{s}}= 3.94\pm 0.13V
\end{equation}
\end{Large}\\

\end{document}